\subsection{R* 中可乘的族}

在有限非零实数的乘法群 \(\mathbb{R}^*\) 中，族 \((x_{t})_{i\in I}\) 只当 \(\lim x_{t} = 1\) （关于 I 的有限子集的余集组成的滤子）时才可能可乘（第三章，§ 5，第 2 节，命题 1）。特别地，指标 \(i\) （使 \(x_{t} < 0\) ）只能有有限多个。因此我们可以限于只考虑各项都严格为正的族 \(x_{t}\) ；这时令 \(x_{i} = i + u_{i}\) 是方便的，其中诸 \(u_{i}\) 满足不等式 \(-i < u_{i} < +\infty\) （对一切 \(i\) ）。由于 \(\mathbf{R}^*\) 的每一点都有可数的基本邻域系，指标 \(i\) （使 \(u_{i} \neq 0\) ）组成的集合是可数的，只要族 \((i + u_{i})\) 在 \(\mathbf{R}^*\) 中可乘。

定理 4. 族 \((1 + u_{t})\) 在 \(\mathbb{R}^*\) 中可乘当且仅当族 \((u_{t})\) 在 \(\mathbb{R}\) 中可和。

引理. (i) 若 \((a_{i})_{1\leqslant i\leqslant p}\) 是数 \(>0\) 的一个有限序列，则

\[
\prod_ {i = 1} ^ {p} \left(1 + a _ {i}\right) \geqslant 1 + \sum_ {i = 1} ^ {p} a _ {i}.\tag{2}
\]

\begin{enumerate}
\def\labelenumi{(\roman{enumi})}
\setcounter{enumi}{1}
\tightlist
\item
  若又有 \(a_i < 1\) 对每个 \(i\) 成立，则
\end{enumerate}

\[
\prod_ {i = 1} ^ {p} \left(\mathrm{I} - a _ {i}\right) \geqslant \mathrm{I} - \sum_ {i = 1} ^ {p} a _ {i}.\tag{3}
\]

这些不等式当 \(p = 1\) 时是显然的，并可对 \(p\) 用归纳法证明。若

\[
\prod_ {i = 1} ^ {p - 1} \left(\mathrm{I} + a _ {i}\right) \geqslant \mathrm{I} + \sum_ {i = 1} ^ {p - 1} a _ {i},
\]

则我们有

\[
\prod_ {i = 1} ^ {p} \left(\mathrm{I} + a _ {i}\right) \geqslant \left(\mathrm{I} + a _ {p}\right) \left(\mathrm{I} + \sum_ {i = 1} ^ {p - 1} a _ {i}\right) = \mathrm{I} + \sum_ {i = 1} ^ {p} a _ {i} + a _ {p}. \sum_ {i = 1} ^ {p - 1} a _ {i} \geqslant \mathrm{I} + \sum_ {i = 1} ^ {p} a _ {i}.
\]

同样，若

\[
\prod_ {i = 1} ^ {p - 1} \left(\mathrm{I} - a _ {i}\right) \geqslant \mathrm{I} - \sum_ {i = 1} ^ {p - 1} a _ {i},
\]

则我们有

\[
\prod_ {i = 1} ^ {p} \left(\mathrm{I} - a _ {i}\right) \geqslant \left(\mathrm{I} - a _ {p}\right) \left(\mathrm{I} - \sum_ {i = 1} ^ {p - 1} a _ {i}\right) = \mathrm{I} - \sum_ {i = 1} ^ {p} a _ {i} + a _ {p}. \sum_ {i = 1} ^ {p - 1} a _ {i} \geqslant \mathrm{I} - \sum_ {i = 1} ^ {p} a _ {i}.
\]

证明了引理之后，我们注意，若族 \((\mathbf{I} + u_{t})\) 可乘，则族 \((\mathbf{I} + u_{t}^{+})\) 与 \((\mathbf{I} - u_{t}^{-})\) 也可乘，因为 \(\mathbf{R}^*\) 是完备群（第三章，§ 5，第 3 节，命题 2）；反之，若族 \((\mathbf{I} + u_{t}^{+})\) 与 \((\mathbf{I} - u_{t}^{-})\) 可乘，则 \((\mathbf{I} + u_{t})\) 也可乘（第三章，§ 5，第 3 节，命题 3）。因此我们只需考虑两种情形：一切 \(u_{t}\) 都 \(\geqslant 0\) ，或它们都 \(\leqslant 0\) 。

首先设 \(u_{t} \geqslant 0\) 对一切 \(t\) 成立。若族 \((1 + u_{t})\) 可乘，则对每个 \(\varepsilon > 0\) ，存在一个有限子集 \(J\) （属于指标集 \(I\) ），使得对每个有限子集 \(H\) （含于 \(I\) 且与 \(J\) 不相交），有 \(\mathbf{I} \leqslant \prod_{i \in \mathbf{H}} (\mathbf{I} + u_i) \leqslant \mathbf{I} + \varepsilon\) ；由 (2) 推出 \(\sum_{i \in \mathbf{H}} u_i \leqslant \varepsilon\) ，这表明由 Cauchy 判别准则（第三章，§ 5，第 2 节，定理 I），\((u_i)\) 在 \(\mathbf{R}\) 中可和。

反之，设 \((u_{t})\) 在 \(\mathbf{R}\) 中可和。对每个 \(\varepsilon\) （满足 \(0 < \varepsilon < 1\) ），存在一个有限子集 \(J\) （\(I\) 的），使得对每个有限子集 \(H\) （含于 \(I\) 且与 \(J\) 不相交），有 \(0 \leqslant \sum_{t \in H} u_{t} \leqslant \varepsilon\) 。由 (3)，我们因此有 \(\prod_{t \in H} (I - u_{t}) \geqslant I - \varepsilon\) ；但由于 \(I + u \leqslant \frac{I}{I - u}\) 对每个数 \(u\) （满足 \(0 \leqslant u < 1\) ）成立，由此推出

\[
\mathrm{I} \leqslant \prod_ {i \in \mathrm{H}} (\mathrm{I} + u _ {i}) \leqslant \frac {\mathrm{I}}{\mathrm{I} - \varepsilon},
\]

而这表明 \((\mathbf{I} + u_{t})\) 可乘（Cauchy 判别准则）。

当一切 \(u_{i}\) 都 \(\leqslant 0\) 时证明类似。为证明 \((u_{i})\) 可和当 \((1 + u_{i})\) 可乘，用公式 (2) 及不等式 \(1 - u \leqslant \frac{1}{1 + u}\) ( \(0 \leqslant u < 1\) )；为证明 \((1 + u_{i})\) 可乘当 \((u_{i})\) 可和，用公式 (3)。

在第五章（§ 4）中，群 \(\mathbf{R}^*\) 的拓扑研究将使我们能借助对数函数给出族在 \(\mathbf{R}^*\) 中可乘的另一个判别准则；在后续分册中我们将利用对数的微分性质证明这个判别准则与上面那个等价。

